\begin{frame}{Digitale Signaturen mit RSA}
	Neben der Verschlüsselung können asymmetrische Verfahren auch für \textbf{digitale Signaturen} genutzt werden.
	\pause
	\begin{mydefinition}[Ziele digitaler Signaturen]
		\pause
		\begin{itemize}[<+->]
			\item \textbf{Authentizität}: Garantiert die Identität der absendenden Person (nur der Besitzer des passenden privaten Schlüssels konnte die Signatur erzeugen).
			\item \textbf{Integrität}: Garantiert, dass die Nachricht auf dem Übertragungsweg nicht verändert oder manipuliert wurde.
			\item \textbf{Verbindlichkeit (Nicht-Abstreitbarkeit)}: Die Urheberschaft kann im Nachhinein nicht geleugnet werden.
		\end{itemize}
	\end{mydefinition}
\end{frame}

\begin{frame}{Funktionsweise von RSA-Signaturen}
	Beim Signieren wird die Rolle der Schlüssel im Vergleich zur Verschlüsselung vertauscht:
	\pause
	\begin{myprocedure}[Digitale Signatur mit RSA (Grundprinzip)]
		{\footnotesize
		\pause
		\begin{enumerate}[<+->]
			\item \textbf{Signaturerstellung (Absenderin Alice)}:
			      Alice signiert den Klartext $m$ mit ihrem \textbf{privaten Schlüssel} $d_A$:
			      \[ s = m^{d_A} \bmod n_A \]
			      Sie sendet die Nachricht $m$ zusammen mit der Signatur $s$ an Bob.
			\item \textbf{Signaturüberprüfung (Empfänger Bob)}:
			      Bob prüft die Signatur mit Alices \textbf{öffentlichem Schlüssel} $(e_A, n_A)$:
			      \[ m' = s^{e_A} \bmod n_A \]
			      Gilt $m' = m$, ist die Signatur mathematisch gültig, unverfälscht und stammt garantiert von Alice!
		\end{enumerate}
		}
	\end{myprocedure}
\end{frame}

\begin{frame}{Beispiel: RSA Digitale Signatur}
	\begin{exampleblock}{Rechenbeispiel (aus dem Skript)}
		{\small
		Alice besitzt den privaten Schlüssel $(d, n) = (9677, 12317)$ und den öffentlichen Schlüssel $(e, n) = (5, 12317)$. Sie möchte die Nachricht $m = 42$ signieren.
		
		\begin{itemize}[<+->]
			\item \textbf{Signaturerstellung (Alice)}:
			      \[ s = m^d \bmod n = 42^{9677} \bmod 12317 = 161 \]
			      Alice übermittelt das Paar $(m, s) = (42, 161)$ an Bob.
			\item \textbf{Signaturüberprüfung durch Bob}:
			      \[ m' = s^e \bmod n = 161^5 \bmod 12317 = 42 \]
			      Da $m' = m = 42$, ist die Signatur authentisch und unverfälscht!
			\item \textbf{Erkennung von Manipulation}:
			      Verändert Angreiferin Eve die Nachricht zu $m = 43$ (mit altem $s=161$), so ergibt die Prüfung $161^5 \bmod 12317 = 42 \neq 43$. Bob erkennt die Fälschung sofort!
		\end{itemize}
		}
	\end{exampleblock}
\end{frame}

\begin{frame}{Auftrag: Digitale Signatur}
	\aufgabebox{4.2 im Skript (RSA Digitale Signatur)}
\end{frame}

\begin{frame}{Praxis-Problem \& Lösung: Kryptographische Hashfunktionen}
	\begin{itemize}[<+->]
		\item \textbf{Problem beim direkten Signieren von $m$}:
		      \begin{enumerate}
			      \item RSA kann nur Zahlen $m < n$ verarbeiten. Reale Dokumente (PDFs, Verträge, E-Mails) sind viele Megabytes gross.
			      \item Modulare Potenzierung mit grossen Exponenten ist rechenintensiv und langsam.
		      \end{enumerate}
		\item \textbf{Lösung: Hashwert signieren}:
		      \begin{itemize}
			      \item Eine \textbf{kryptographische Hashfunktion} (z.\,B.\ \textbf{SHA-256}) komprimiert ein beliebig grosses Dokument auf einen kurzen, eindeutigen Fingerabdruck fester Länge (z.\,B.\ 256 Bit).
			      \item \textbf{Eigenschaften}: Einwegfunktion (nicht umkehrbar), kollisionsresistent, Lawineneffekt.
		      \end{itemize}
		\item \textbf{Ablauf in der Praxis}:
		      \begin{align*}
		          \text{Alice signiert:} &\quad s := H(m)^{d_A} \bmod n_A \\
		          \text{Bob prüft:} &\quad s^{e_A} \bmod n_A \stackrel{?}{=} H(m)
		      \end{align*}
	\end{itemize}
\end{frame}

\begin{frame}{Vertraulichkeit vs. Authentizität: Sign-then-Encrypt}
	\begin{columns}[T]
		\begin{column}{0.48\textwidth}
			\begin{block}{Nur Signatur}
				{\small
				\begin{itemize}
					\item Alice sendet Klartext $m$ + Signatur $s$.
					\item \textbf{Authentizität:} Bob weiss, dass $m$ von Alice stammt.
					\item \textbf{Integrität:} $m$ wurde nicht verändert.
					\item \textbf{Keine Vertraulichkeit:} Jede Person auf dem Kanal kann $m$ lesen.
				\end{itemize}
				}
			\end{block}
		\end{column}
		\begin{column}{0.48\textwidth}
			\begin{block}{Sign-then-Encrypt}
				{\small
				\begin{itemize}
					\item 1. Alice signiert $m$ mit $d_A \rightarrow s$.
					\item 2. Alice verschlüsselt $(m, s)$ mit Bobs Schlüssel $e_B \rightarrow (c_m, c_s)$.
					\item 3. Bob entschlüsselt mit $d_B$.
					\item 4. Bob verifiziert Signatur mit $e_A$.
					\item $\rightarrow$ \textbf{Vertraulich \& Authentisch!}
				\end{itemize}
				}
			\end{block}
		\end{column}
	\end{columns}
\end{frame}

\begin{frame}{Ablauf: Sign-then-Encrypt}
	\begin{figure}[H]
		\centering
		\resizebox{!}{0.75\textheight}{\begin{tikzpicture}[
	scale=0.80, every node/.style={transform shape},
	actor/.style={rectangle, rounded corners, draw=black, fill=gray!15, minimum width=3.4cm, minimum height=1.1cm, font=\bfseries, align=center, inner sep=4pt},
	stepbox/.style={rectangle, rounded corners, draw=#1, fill=#1!8, line width=0.8pt, font=\small, align=center, inner sep=6pt},
	msgarrow/.style={->, >=stealth, line width=1pt, #1},
	badge/.style={circle, fill=#1, text=white, font=\bfseries\footnotesize, inner sep=2pt}
]
	% Spaltenkoordinaten
	\def\colA{0}   % Alice
	\def\colB{9.5} % Bob

	% --- AKTEURE (Kopfzeile) ---
	\node[actor] (alice) at (\colA, 0) {Alice\\{\footnotesize (Absenderin)}\\[2pt]{\scriptsize priv.: $d_A$, Bobs öff.: $(e_B, n_B)$}};
	\node[actor] (bob) at (\colB, 0) {Bob\\{\footnotesize (Empfänger)}\\[2pt]{\scriptsize priv.: $d_B$, Alices öff.: $(e_A, n_A)$}};

	% Vertikale Hilfslinien
	\draw[dashed, gray!40] (alice.south) -- (\colA, -10.2);
	\draw[dashed, gray!40] (bob.south) -- (\colB, -10.2);

	% --- SCHRITT 1 & 2 BEI ALICE ---
	\node[font=\footnotesize\bfseries\color{RoyalBlue!80!black}, fill=white, inner sep=2pt] at (\colA, -1.1) {Phase 1: Signieren \& Verschlüsseln};

	\node[stepbox=RoyalBlue] (alice_proc) at (\colA, -2.7) {
		\textbf{1. Signatur erzeugen (mit $d_A$):}\\
		$s := m^{d_A} \bmod n_A$\\[0.4em]
		\textbf{2. Verschlüsseln (mit $e_B$):}\\
		$c_m := m^{e_B} \bmod n_B$\\
		$c_s := s^{e_B} \bmod n_B$
	};

	% --- NACHRICHTENÜBERTRAGUNG ---
	\node[badge=DarkOrchid] (b_msg) at (\colA+0.6, -5.1) {\faLock};
	\draw[msgarrow=DarkOrchid] (b_msg.east) -- (\colB-0.6, -5.1)
		node[midway, above=3pt, font=\footnotesize\bfseries] {Chiffrate $(c_m, c_s)$}
		node[midway, below=3pt, font=\scriptsize\color{gray!70!black}] {(vertraulich \& verschlüsselt: Eve kann weder $m$ noch $s$ lesen)};

	% --- SCHRITT 3 & 4 BEI BOB ---
	\node[font=\footnotesize\bfseries\color{ForestGreen!80!black}, fill=white, inner sep=2pt] at (\colB, -6.8) {Phase 2: Entschlüsseln \& Verifizieren};

	\node[stepbox=ForestGreen] (bob_proc) at (\colB, -8.6) {
		\textbf{3. Entschlüsseln (mit $d_B$):}\\
		$m := (c_m)^{d_B} \bmod n_B$\\
		$s := (c_s)^{d_B} \bmod n_B$\\[0.4em]
		\textbf{4. Signatur prüfen (mit $e_A$):}\\
		$m' := s^{e_A} \bmod n_A$\\
		$\rightarrow \text{Gilt } m' = m\text{?} \implies \text{Echt, unverfälscht \& geheim!} \textcolor{ForestGreen}{\mathbf{\checkmark}}$
	};
\end{tikzpicture}
}
	\end{figure}
\end{frame}

\begin{frame}{Auftrag: Sign-then-Encrypt}
	\aufgabebox{Aufgabe 4.3 im Skript (Sign-then-Encrypt)}
\end{frame}
